\documentclass[10pt,a4paper,twocolumn]{ctexart}
\usepackage[a4paper,left=1.55cm,right=1.55cm,top=1.75cm,bottom=1.75cm,columnsep=0.65cm]{geometry}
\usepackage{amsmath,amssymb}
\usepackage{booktabs,tabularx,array}
\usepackage{caption}
\usepackage{cite}
\usepackage{titlesec}
\usepackage[hidelinks]{hyperref}
\usepackage{stfloats}
\setlength{\parindent}{2em}
\setlength{\parskip}{0pt}
\linespread{0.98}
\raggedbottom
\renewcommand{\tablename}{表}
\renewcommand{\thetable}{\Roman{table}}
\renewcommand{\figurename}{图}
\captionsetup[table]{font=small,labelfont=bf,labelsep=space}
\titleformat{\section}{\centering\normalsize\bfseries}{\Roman{section}.}{0.5em}{}
\titleformat{\subsection}{\normalsize\bfseries}{\Alph{subsection}.}{0.5em}{}
\titlespacing*{\section}{0pt}{0.9em}{0.45em}
\titlespacing*{\subsection}{0pt}{0.65em}{0.25em}
\setlength{\textfloatsep}{0.8em}
\setlength{\floatsep}{0.7em}
\setlength{\intextsep}{0.7em}
\setlength{\abovecaptionskip}{0.25em}
\setlength{\belowcaptionskip}{0.25em}
\begin{document}
\twocolumn[
\begin{@twocolumnfalse}
\begin{center}
{\LARGE\bfseries DIAL：面向边缘异构平台的多模态大模型分布式协同推理框架\par}
\vspace{0.75em}
\begin{tabular*}{\textwidth}{@{\extracolsep{\fill}}ccc@{}}
\parbox{0.31\textwidth}{\centering\normalsize 作者A\\ \small 单位A\\ \small 邮箱A} & \parbox{0.31\textwidth}{\centering\normalsize 作者B\\ \small 单位B\\ \small 邮箱B} & \parbox{0.31\textwidth}{\centering\normalsize 作者C\\ \small 单位C\\ \small 邮箱C} \\
\end{tabular*}
\vspace{0.45em}
\end{center}
\vspace{-0.35em}
\noindent\textbf{摘要—}边缘多媒体应用正在从固定类别感知走向开放语义理解，视觉语言模型为智能监控、工业巡检和移动机器人提供了图像理解、视觉问答和自然语言交互能力。然而，以 Qwen3-VL 为代表的多模态大模型同时包含视觉编码器、语言解码器、跨模态特征融合和自回归缓存等组件，在边缘异构设备上部署时面临模型权重占用高、首 token 延迟长、硬件后端能力不均衡以及跨节点通信开销显著等问题。单设备执行难以同时满足内存容量和实时响应要求，而直接分布式执行又会因中间特征传输和远程调用频繁而降低收益。为此，本文提出 DIAL，一种面向边缘异构平台的多模态大模型分布式协同推理框架。DIAL 将边缘多模态推理抽象为“层段预设、层段聚合、异构辅助”的三阶段协同问题：首先在部署阶段根据设备资源确定连续 Transformer 层段归属，避免运行时动态切分模型带来的后端适配复杂性；其次在推理阶段将同一远端节点上的连续层段合并为一次协同执行，以降低自回归解码中的通信往返次数；最后将视觉编码器和规则文本子图映射到边缘 NPU，并以软件后端完成动态缓存、控制流和回退执行。针对多模态输入，DIAL 进一步设计视觉 token 区间对齐、深层视觉特征注入和首 token 延迟约束下的图像尺度选择机制，使视觉特征融合、输入长度控制和硬件执行路径形成统一优化目标。本文从层段划分、通信开销、多模态融合、视觉尺度控制和异构执行五个方面展开设计，并给出面向真实边缘设备的实验评估方案。DIAL 为多模态大模型在边缘异构环境中的低成本部署提供了一种可运行、可观测和可扩展的系统路径。

\noindent\textbf{关键词—}边缘计算；多模态大模型；分布式推理；异构计算；Qwen3-VL；RKNN；视觉语言模型
\vspace{1em}
\end{@twocolumnfalse}
]

\section{引言}

智能监控、工业巡检和移动机器人等边缘多媒体应用正在从传统感知任务扩展到开放语义理解任务。传统视觉系统通常输出检测框、类别标签或异常标志，而多模态大模型能够根据图像内容直接回答自然语言问题。例如，在安防场景中，系统需要判断画面中是否存在摔倒行为并给出解释；在工业场景中，系统需要根据巡检图像描述设备异常；在移动机器人场景中，系统需要结合视觉输入完成环境理解和任务交互。这类任务要求边缘设备不仅具备视觉感知能力，还需要具备语言理解和生成能力。

多模态大模型的边缘部署难度显著高于传统 CNN 模型。首先，模型结构更加复杂。Qwen3-VL 等视觉语言模型由视觉编码器、语言解码器、视觉投影模块和多层跨模态融合结构组成，模型参数和中间状态远大于轻量视觉模型。其次，首 token 延迟较高。图文请求在生成第一个 token 之前，需要经历图片解码、尺寸调整、视觉编码、视觉 embedding 注入以及文本 prefill 等步骤。对于边缘设备而言，首 token 延迟直接影响交互体验。再次，硬件后端能力存在差异。CPU 具备通用性但吞吐有限，GPU 适合矩阵计算但边缘平台算力有限，NPU 在特定静态子图上具备能效优势，但受输入形状、算子支持和模型格式限制，难以直接承载完整语言模型。

分布式推理可以缓解单设备内存和算力不足问题。通过将 Transformer 层划分到多个节点执行，系统能够降低单个设备的模型加载压力，并利用远端设备参与计算。然而，边缘分布式推理并不等价于云端 GPU 集群并行。边缘节点之间通常通过以太网或无线链路连接，带宽、时延和稳定性有限；Transformer 的中间 hidden state 需要跨节点传输，如果每一层都触发远程调用，通信开销将在自回归解码阶段被反复放大。此外，部分边缘推理后端以编译后的模型文件运行，运行时任意切分模型或动态迁移层并不可行。

本文提出 DIAL，目标是在不改变模型结构和不重新训练模型的前提下，实现多模态大模型在边缘异构平台上的协同推理。DIAL 的核心思想是离线预设层划分和在线轻量协同执行。离线阶段根据设备资源和部署需求确定每个节点负责的层段；在线阶段，主控节点按照层顺序执行模型，当遇到远端连续层段时，将中间 hidden state 发送给对应执行节点，由执行节点连续完成该层段计算后返回结果。该方式避免了运行时动态切分模型的复杂性，同时减少了远程调用次数。

除分布式层执行外，DIAL 还针对多模态模型和异构硬件进行了设计。对于 Qwen3-VL 图文输入，DIAL 通过视觉 token 区间记录和 embedding 注入，使视觉编码器输出能够准确写入语言模型输入序列；对于 deepstack 等深层视觉特征，DIAL 在语言模型对应层中进行二次注入，保持原模型跨模态融合方式。对于 RK3588 等边缘异构平台，DIAL 将视觉编码器和部分文本子图映射到 RKNN Runtime 上执行，同时保留软件路径作为回退机制，从而在性能和稳定性之间取得平衡。

本文贡献如下。

\begin{itemize}
\setlength{\itemsep}{0.12em}
\setlength{\parsep}{0pt}
\setlength{\topsep}{0.25em}
\item 提出一种“层段预设、层段聚合、异构辅助”的三阶段协同推理框架，将边缘多模态大模型部署问题表述为层段归属、跨节点通信和异构后端选择的联合优化问题。
\item 设计资源约束下的预设层段划分机制，将运行时模型动态切分问题转化为部署阶段拓扑配置问题，在保持 Transformer 顺序语义不变的前提下，将连续层段部署到主控节点和多个执行节点。
\item 提出通信开销感知的层段聚合机制，根据远端层段的节点归属合并相邻层执行，将解码阶段的远程调用次数由与远程层数相关转化为与连续层段数量相关，从而降低通信往返和协议处理开销。
\item 提出面向首 token 延迟的多模态输入控制机制，在视觉 token 区间对齐和 deepstack 特征融合基础上，引入图像尺度约束策略，使视觉编码耗时、语言模型 prefill 长度和任务语义质量能够在边缘设备上进行可控折中。
\item 提出 NPU 辅助的混合执行策略，将视觉编码器、文本连续层段和部分文本子图接入 RKNN 后端，在边缘 NPU 可用时利用硬件加速，在后端不适配时回退到软件路径，增强系统在真实异构平台上的部署稳定性。
\end{itemize}

本文其余部分组织如下。第二节介绍边缘分布式推理、大模型推理系统、异构硬件加速和多模态模型边缘部署等相关工作。第三节给出 DIAL 的问题定义、总体框架和关键机制。第四节设计实验方案和评价指标。第五节讨论方法适用范围与局限性。第六节总结全文并展望后续工作。

\section{相关工作}

\subsection{边缘智能与分布式推理}

边缘智能通过在靠近数据源的位置执行模型推理，降低云端传输延迟并减少隐私风险。早期研究主要面向 CNN 分类、检测和分割模型，采用模型压缩、输入自适应、早退推理或端云分割等方法降低端侧负载。Neurosurgeon 等工作将神经网络划分为端侧部分和云端部分，端侧输出中间特征后传输给远端继续计算 \cite{ref1}。后续研究进一步讨论了多节点协同推理和中间特征压缩问题 \cite{ref2}, \cite{ref3}。该类方法在降低端侧计算方面有效，但中间特征传输可能成为新的瓶颈。

分布式推理的划分粒度包括算子级、层级和模块级。算子级划分具有更大搜索空间，但需要复杂运行时和调度器；模块级划分实现简单，但对性能优化空间有限；层级划分在部署灵活性和实现复杂度之间更适合边缘环境。Transformer 结构由多个顺序堆叠的层组成，适合按层划分到不同设备。与面向通用 DNN 的端云切分不同，DIAL 关注的是视觉语言模型在边缘设备上的层级协同执行，并进一步将连续层段作为远程执行单元，以降低自回归解码过程中的通信频率。

\subsection{大模型推理系统}

大模型推理系统通常围绕显存管理、KV cache、批处理调度和算子优化展开。以 PagedAttention 为代表的服务系统通过改进 KV cache 管理提升服务器端大模型服务吞吐 \cite{ref5}，TensorRT-LLM 等框架则针对 GPU 推理图优化和高性能算子执行进行加速 \cite{ref6}。这类系统主要面向高性能 GPU，强调高并发吞吐和显存利用率。然而，边缘设备的计算能力、内存容量和网络条件与服务器环境差异明显。边缘推理更关注单请求延迟、部署稳定性和硬件后端适配。特别是多模态模型推理不仅包括语言模型解码，还包括视觉编码和图文融合，首 token 延迟受到图像尺寸、视觉 token 数量和视觉后端性能共同影响。

现有大模型服务框架大多假设模型运行在单个或多个高速互联 GPU 上，对边缘 NPU、低功耗 CPU 和跨设备中间特征传输关注不足。部分边缘 AI 部署研究分析了黑盒部署策略对延迟和模型性能的影响 \cite{ref4}，但对视觉语言模型中的视觉 token、首 token 延迟和自回归层段协同仍缺少针对性设计。因此，在边缘环境中部署多模态大模型，需要同时考虑模型层划分、通信开销和硬件后端兼容性。

\subsection{异构硬件推理加速}

边缘 SoC 通常集成 CPU、GPU 和 NPU。NPU 在固定形状和固定算子图上具有较高能效，但其部署依赖专用工具链，例如 RKNN Toolkit 和 RKNN Runtime \cite{ref8}。完整大语言模型直接迁移到 NPU 面临多方面困难：自回归解码导致上下文长度动态变化，KV cache 需要持续更新，注意力 mask 和 RoPE 等机制对后端支持要求较高，完整模型体积也可能超过 NPU 可承载范围。

因此，更可行的方式是将局部计算密集子图映射到 NPU，而将控制流、动态缓存和不稳定算子保留在软件后端。视觉编码器通常具有较稳定输入形状，是较适合 NPU 加速的模块；语言模型中的 QKV 投影和 MLP 子图也具有较强矩阵计算特征，可作为局部加速对象。DIAL 采用这种混合执行策略，以软件路径保证正确性和稳定性，以 NPU 子图降低关键路径耗时。与单纯追求全模型 NPU 化不同，DIAL 将异构执行视为边缘协同推理框架中的一部分，使 NPU 加速、层段预设和软件回退能够共同服务于端到端延迟目标。

\subsection{多模态模型边缘部署}

多模态大模型的边缘部署不仅需要执行语言模型解码，还需要处理图片解码、尺度调整、视觉编码、视觉 token 插入和跨模态特征融合。Qwen-VL 等视觉语言模型表明，将视觉特征以 token 或 embedding 形式注入语言模型是实现开放式图文理解的重要路径 \cite{ref7}。与纯文本大语言模型相比，视觉语言模型的首 token 延迟更容易受到输入图像尺度和视觉 token 数量影响。现有部署实践通常采用固定输入分辨率以简化工程实现，但固定尺度难以同时适配监控告警、巡检问答和机器人交互等不同延迟需求。

因此，边缘多模态推理需要同时保证模型语义一致性和输入计算量可控。视觉特征必须按照原模型结构写入语言模型输入序列，避免退化为文本描述式图像理解；同时，系统需要根据硬件后端和任务响应要求控制视觉 token 数量。DIAL 围绕这一问题设计视觉 token 区间对齐、deepstack 特征融合和首 token 延迟约束下的图像尺度控制，使多模态输入处理成为边缘推理优化的一部分。

\section{框架与设计}

\subsection{问题定义与符号说明}

本文研究边缘异构环境下的多模态大模型协同推理问题。给定一个视觉语言模型 \(F\)，其推理过程包括视觉编码、跨模态特征融合、语言模型 prefill 和自回归 decode。边缘设备集合记为 \(D=\{d_0,d_1,\ldots,d_n\}\)，其中 \(d_0\) 为主控节点，其他设备为可选执行节点。语言模型包含 \(L\) 个顺序 Transformer 层，记为 \(M=\{l_0,l_1,\ldots,l_{L-1}\}\)。对于输入请求 \(x=(x_t,x_v)\)，其中 \(x_t\) 表示文本输入，\(x_v\) 表示图像输入，系统需要在设备资源、网络链路和硬件后端限制下完成端到端生成。

DIAL 的目标是在保持原模型参数和计算语义不变的前提下，降低端到端延迟和分布式非计算开销。端到端延迟可表示为：

\[
T_{e2e}=T_{prep}+T_{vis}+T_{prefill}+\sum_{i=1}^{N}T_{decode}^{(i)},
\]
其中 \(T_{prep}\) 表示输入预处理耗时，\(T_{vis}\) 表示视觉编码耗时，\(T_{prefill}\) 表示首 token 前的语言模型预填充耗时，\(T_{decode}^{(i)}\) 表示第 \(i\) 个生成 token 的解码耗时，\(N\) 为生成 token 数量。对于边缘交互式任务，首 token 延迟 \(T_{first}=T_{prep}+T_{vis}+T_{prefill}+T_{decode}^{(1)}\) 是关键指标。

该问题受到三类约束。第一是资源约束，分配到设备 \(d_k\) 的层段权重、KV cache 和运行时缓冲不能超过设备可用内存 \(C_k\)。第二是通信约束，跨节点层段执行会引入序列化、网络传输和反序列化开销，若层段切分过细，远端计算收益会被通信往返抵消。第三是后端约束，NPU 等编译式后端通常要求固定输入形状和受支持算子集合，不能任意承载完整动态语言模型。因此，DIAL 不求解全局动态调度最优解，而是面向静态或半静态边缘部署，构建可配置、可测量和可回退的协同推理策略。

表 I 给出本文主要符号说明。

\begin{table}[t]
\centering
\caption{主要符号说明}
\small
\begin{tabular}{ll}
\toprule
符号 & 含义 \\
\midrule
\(F\) & 待部署的视觉语言模型 \\
\(D\) & 边缘设备集合 \\
\(d_0\) & 主控节点 \\
\(M\) & 语言模型 Transformer 层集合 \\
\(L\) & Transformer 层数 \\
\([s,e]\) & 起始层为 \(s\)、结束层为 \(e\) 的连续层段 \\
\(C_k\) & 设备 \(d_k\) 的可用内存 \\
\(W_{s:e}\) & 层段 \([s,e]\) 的权重和缓存占用 \\
\(S\) & 候选图像尺度集合 \\
\(T_{first}\) & 首 token 延迟 \\
\(T_{dist}\) & 分布式非计算开销 \\
\(T_{max}\) & 应用允许的首 token 延迟上界 \\
\bottomrule
\end{tabular}
\end{table}

\subsection{DIAL 总体方法框架}

DIAL 在边缘异构约束下重新组织多模态大模型前向过程。给定视觉语言模型 \(F\)、边缘设备集合 \(D\)、网络链路状态 \(B\) 和输入请求 \(x=(x_t,x_v)\)，DIAL 需要在不改变模型参数和计算语义的前提下，确定三个关键决策：语言模型层段由哪些设备执行，跨节点执行边界如何合并，以及视觉和文本子图应选择何种硬件后端。

为此，DIAL 将协同推理过程抽象为三阶段框架。第一阶段为层段预设。系统在部署阶段根据设备内存、计算能力和后端可用性，将 Transformer 层划分为若干连续层段，并为每个层段指定执行节点。第二阶段为层段聚合。在线推理时，主控节点沿模型层顺序推进计算，遇到远端连续层段时只触发一次协同执行，由远端节点完成该层段内的全部计算后返回输出。第三阶段为异构辅助。系统将输入形状稳定、算子结构规则的视觉编码器和部分文本子图映射到 NPU，将动态缓存、控制流和不稳定算子保留在 CPU/GPU 软件路径。

该框架的核心在于将边缘推理中的三个矛盾统一到同一执行模型中。模型规模与单机资源之间的矛盾由层段预设缓解；分布式计算收益与通信代价之间的矛盾由层段聚合缓解；硬件加速能力与后端适配限制之间的矛盾由异构辅助缓解。与只强调系统部署流程的方案不同，DIAL 的方法主线围绕层段、张量和硬件后端三个对象展开，使后续优化能够落在明确的可测指标上，包括首 token 延迟、decode 吞吐、远端计算耗时、通信非计算开销和 NPU 覆盖比例。

\subsection{资源约束下的层段预设}

DIAL 面向的边缘环境通常具有两个约束：一是单节点无法稳定加载完整多模态模型，二是部分硬件后端不支持运行时任意修改模型结构。因此，DIAL 不采用运行时动态切分模型，而是在部署阶段确定层划分方案。设语言模型包含 \(L\) 个 Transformer 层，记为 \(M=\{l_0,l_1,\ldots,l_{L-1}\}\)，边缘执行节点集合为 \(D=\{d_0,d_1,\ldots,d_n\}\)。预设划分将模型层映射到不同节点，其中节点 \(d_k\) 负责执行连续层段 \(\{l_s,l_{s+1},\ldots,l_e\}\)。

连续层段是 DIAL 的基本部署单元。选择连续层段而非离散层集合有两个原因。首先，Transformer 前向过程具有强顺序性，连续层段不改变模型执行顺序。其次，连续层段能够让远端节点在一次通信后完成多个相邻层计算，避免中间 hidden state 在节点之间反复往返。部署完成后，每个节点只加载其负责的模型层，从而降低单节点内存占用。

在该设计下，层划分从运行时调度问题转化为部署阶段配置问题。对于资源较弱的主控节点，可以将前若干层或中间层部署到远端节点；对于仅单机运行的场景，可以将所有层保留在本地。该策略牺牲了部分运行时自适应能力，但换来了更高的工程可控性和后端兼容性，适合实际边缘设备中模型文件较大、硬件运行时能力有限的场景。

进一步地，DIAL 在层段选择时遵循资源约束和通信约束。设设备 \(d_k\) 的可用内存为 \(C_k\)，层段 \([s,e]\) 的权重和缓存占用为 \(W_{s:e}\)，则可部署层段需要满足 \(W_{s:e}\le C_k\)。同时，层段长度不宜过短，否则远端计算收益会被固定通信开销抵消；层段长度也不宜过长，否则单个执行节点可能成为新的瓶颈。因此，DIAL 优先选择能够被目标设备稳定加载、且在解码阶段具有足够计算量的连续层段。该原则为后续实验中的不同层段划分对比提供依据。

层段预设过程包括四个步骤。首先，系统记录各设备的可用内存、计算后端和可承载模型格式，形成设备能力描述。其次，根据模型层数、单层权重规模和 KV cache 需求估计不同连续层段的资源占用。再次，筛选满足设备内存约束和后端约束的候选层段，排除过短且通信收益较低的切分方案。最后，根据主控节点负载、执行节点能力和网络条件确定最终层段归属。该过程既可以由人工根据 profiling 结果配置，也可以作为后续自动层划分算法的输入。

\subsection{通信开销感知的层段聚合}

自回归语言模型的 decode 阶段具有重复执行特征：每生成一个 token，都需要按照顺序执行所有 Transformer 层。若远端节点负责 \(m\) 个层且每层单独触发一次远程调用，则生成 \(N\) 个 token 需要近似 \(N\cdot m\) 次通信。通信往返将成为显著开销，尤其在边缘网络带宽和时延不稳定时更为明显。

DIAL 将同一远端节点上的连续层段合并为一次远程执行。设远端节点负责层段 \([s,e]\)，长度为 \(m=e-s+1\)。主控节点在进入层 \(s\) 前发送当前 hidden state，远端节点连续计算 \(l_s\) 至 \(l_e\)，并返回层 \(e\) 的输出。这样，生成 \(N\) 个 token 的远程调用次数由 \(N\cdot m\) 降低为 \(N\)。该优化不改变模型数学计算，只改变跨节点执行边界。

层段聚合还降低了中间状态传输量。若逐层远程执行，则每一层输出都需要回到主控节点再发送到远端，层间 hidden state 被多次传输；层段聚合执行则将远端层之间的中间结果保留在执行节点本地，仅在层段入口和出口发生跨节点传输。对于带 KV cache 的模型，远端节点只维护自身负责层的缓存，主控节点维护本地层缓存，从而避免跨节点同步完整 KV cache。

从优化角度看，层段聚合可以写成对远程调用边界的压缩。设远程执行层集合为 \(R\)，相邻层 \(l_i\) 和 \(l_{i+1}\) 属于同一远端节点时，DIAL 将它们合并到同一执行段；当相邻层位于不同节点或从本地切换到远端时，才产生新的跨节点边界。因此，单个 token 的远程调用次数近似为远程连续层段数量，而不是远程层数。对于边缘链路中常见的毫秒级往返时延，这一变化直接影响 decode 阶段的逐 token 响应速度。

在线协同推理过程按照模型层顺序执行。对于本地层，主控节点直接完成前向计算并更新本地 KV cache；对于远端层段，主控节点在层段入口处发送当前 hidden state、位置索引和缓存状态所需的最小信息，执行节点连续完成该层段计算并维护自身层段的 KV cache，然后返回层段出口 hidden state。prefill 阶段和 decode 阶段采用相同的层段边界，但 decode 阶段会在每个生成 token 上重复触发，因此层段聚合对 decode 阶段收益更明显。

\subsection{通信开销建模与轻量传输}

边缘分布式推理的关键瓶颈之一是中间特征传输。语言模型层间传输的是高维 hidden state，而不是文本 token。若采用文本化协议或冗余字段编码，会增加序列化和解析成本。DIAL 采用面向张量恢复的轻量二进制传输方式，每次远程调用只保留数据缓冲区、数据类型和维度信息，使接收端能够在目标设备上恢复中间张量。

远程调用总耗时可近似表示为：

\[
T_{remote}=T_{ser}+T_{send}+T_{compute}+T_{recv}+T_{deser},
\]
其中 \(T_{ser}\) 和 \(T_{deser}\) 表示序列化与反序列化耗时，\(T_{send}\) 和 \(T_{recv}\) 表示网络发送与接收耗时，\(T_{compute}\) 表示远端模型计算耗时。为了评估通信和协议部分的影响，DIAL 定义分布式非计算开销：

\[
T_{dist}=T_{remote}-T_{compute}.
\]
该指标用于衡量网络传输、协议处理和张量编解码造成的额外开销。当 \(T_{dist}\) 占比较高时，说明系统瓶颈主要来自跨节点通信；当 \(T_{compute}\) 占比较高时，说明远端计算仍是主要耗时来源。

DIAL 的通信优化主要体现在两个方面。第一，层段聚合减少远程调用次数，从而减少多次发送、接收和协议处理开销。第二，紧凑层段描述减少请求元数据，尤其在 decode 阶段单次输入张量较小时，固定元数据开销更容易影响整体延迟。两者结合，使分布式执行更适合自回归生成过程。

需要指出的是，DIAL 的轻量传输只改变中间特征的表示和传输方式，不改变 hidden state 的数值语义。接收端根据数据类型、维度和连续数据缓冲恢复张量，并在目标设备上继续执行对应层段。该设计避免了将中间特征转换为文本或冗余结构化字段，从而减少协议开销和解析开销。

\subsection{多模态视觉特征融合}

Qwen3-VL 的多模态推理不能简单地将图片转换为文本描述，而需要将视觉编码器输出的 embedding 注入语言模型输入序列。DIAL 为每张图片在文本序列中保留连续视觉 token 区间，并记录该区间的位置。视觉编码器生成的主视觉 embedding 被写入该区间，使语言模型在 prefill 阶段能够直接处理视觉特征。

对于图片输入，DIAL 首先将图片解码为 RGB 图像，并根据后端选择合适尺寸策略。软件视觉路径可采用最大边长限制，以降低视觉 token 数量；NPU 视觉路径采用固定边长补边，以满足静态输入形状约束。视觉编码器输出包括主视觉 embedding 和 deepstack 中间特征。主视觉 embedding 注入输入序列；deepstack 特征在语言模型指定层注入，从而保持 Qwen3-VL 原有的深层视觉融合路径。

形式化地，设语言模型输入 embedding 为 \(H^0\)，第 \(j\) 张图片在输入序列中对应的视觉 token 区间为 \([a_j,b_j]\)，视觉编码器输出的主视觉 embedding 为 \(V_j\)。DIAL 在 prefill 前执行区间对齐，将 \(H^0_{a_j:b_j}\) 替换为 \(V_j\)。对于 deepstack 特征，设其在语言模型层集合 \(\mathcal{R}\) 中进行注入，则在层 \(r\in\mathcal{R}\) 处将对应视觉区间的隐藏状态与深层视觉特征对齐融合。该过程保持了文本 token 与视觉 token 在同一序列中的相对位置关系，也避免了视觉特征在分布式层执行时丢失。

该设计使图文融合与模型结构保持一致。视觉信息不是作为额外文本提示参与生成，而是以 embedding 形式进入语言模型隐藏空间；deepstack 特征进一步增强了视觉信息在语言模型中后层的表达能力。对于边缘多模态任务，这种方式能够在保持模型能力的同时，为视觉编码和文本解码分别选择不同硬件后端。

\subsection{首 token 延迟约束的图像尺度控制}

图文请求的首 token 延迟主要由图片预处理、视觉编码、视觉特征注入和文本 prefill 共同决定。其中，图像尺度会同时影响视觉编码耗时和视觉 token 数量：尺度越大，视觉细节越充分，但视觉编码和语言模型 prefill 的计算量也越高；尺度越小，首 token 延迟更低，但小目标、文字和细节理解可能受损。因此，DIAL 将图像尺度控制作为多模态边缘推理中的独立优化机制，而不是简单的预处理参数。

DIAL 采用首 token 延迟约束下的候选尺度选择策略。设候选图像边长集合为 \(S=\{s_1,s_2,\ldots,s_q\}\)，给定应用允许的首 token 延迟上界 \(T_{max}\)，系统通过离线或启动阶段 profiling 获得不同尺度下的视觉编码耗时 \(T_{vis}(s)\) 和 prefill 估计耗时 \(T_{pre}(s)\)。在线执行时，系统选择满足 \(T_{vis}(s)+T_{pre}(s)\le T_{max}\) 的最大尺度，以尽可能保留视觉细节；若所有候选尺度均超过延迟约束，则选择最小尺度保证响应。

该机制的创新点在于将多模态输入控制与推理关键路径直接关联。传统做法通常给定固定图片尺寸，无法反映不同边缘硬件和不同任务延迟目标的差异；DIAL 则把视觉尺度作为首 token 延迟的可控变量，使系统可以在监控告警、巡检问答和机器人交互等场景中按需选择不同延迟质量折中。对于 NPU 视觉路径，尺度集合受固定输入形状约束；对于软件视觉路径，尺度集合可以更灵活，但仍需满足视觉 token 数量和 prefill 延迟约束。

图像尺度控制的执行过程包括三步。第一，离线测量候选尺度在不同视觉后端上的视觉编码耗时，并估计其对应的视觉 token 数量和 prefill 代价。第二，根据应用场景设置首 token 延迟上界，例如告警类任务可采用更严格的延迟约束，巡检问答类任务可保留更高图像尺度。第三，在线请求到达时选择满足延迟约束的最大候选尺度，并将该尺度传递给视觉编码路径。该策略保留了任务侧对质量和延迟的选择权，避免系统固定在单一图片尺寸上。

\subsection{NPU 辅助异构执行}

DIAL 不将完整模型强行映射到 NPU，而是采用 NPU 辅助的混合执行策略。该策略基于两个观察：第一，视觉编码器通常具有固定输入形状和稳定算子结构，适合转换为 RKNN 模型；第二，语言模型中部分矩阵计算子图可由 NPU 加速，但完整自回归 decode 的动态缓存和注意力机制对 NPU 后端要求较高。

因此，DIAL 将视觉编码器作为优先加速对象。当图文请求进入系统时，视觉编码器可在 NPU 上完成前向计算，输出视觉 embedding 和 deepstack 特征，从而降低首 token 路径中的视觉耗时。对于文本分支，DIAL 支持将连续文本层段或 QKV、MLP 等局部子图映射到 NPU 后端。若当前输入形状、缓存长度或后端能力不满足 NPU 执行要求，系统回退到 CPU/GPU 软件路径。

这种混合策略避免了单一硬件后端失败导致整个推理链路不可用。NPU 负责适合其执行的规则子图，CPU/GPU 负责动态控制、缓存管理和不稳定算子。对于边缘设备而言，该策略比全量 NPU 化更容易落地，也便于根据模型转换结果逐步扩大加速覆盖范围。

异构辅助执行遵循“优先加速、条件校验、失败回退”的原则。系统优先选择视觉编码器等静态子图进行 NPU 加速；在文本分支中，仅当输入长度、缓存形状和算子集合满足目标后端要求时才启用 NPU 子图；若后端初始化失败、输入形状不匹配或运行结果不可用，则自动回退到软件路径。该原则使 DIAL 的异构执行不依赖单一后端成功率，从而更适合真实边缘设备中工具链版本和算子支持不完全稳定的部署环境。

\section{实验设计}

\subsection{实验环境}

实验面向真实边缘异构平台。主控节点运行 DIAL 推理服务并负责对外响应请求，执行节点根据预设层划分加载部分 Transformer 层。边缘设备采用 RK3588 类开发板，NPU 后端使用 RKNN Runtime。模型选择 Qwen3-VL-2B-Instruct 和 Qwen3-VL-8B-Instruct，任务包括纯文本问答和图文问答。网络环境通过以太网连接，并使用 Linux TC 工具限制带宽，以模拟不同边缘链路条件。

实验从五个方面评价系统性能。第一，端到端响应能力，包括首 token 延迟和请求总耗时；第二，生成能力，包括平均 token 生成速率和 decode 阶段生成速率；第三，分布式执行开销，包括远端计算耗时、非计算通信开销、远程调用次数和读写数据量；第四，核心机制的增量收益，包括层段聚合、轻量传输、图像尺度控制和 NPU 辅助执行的消融结果；第五，不同网络带宽下的适用边界。

评价指标定义如下。首 token 延迟表示请求发出到第一个生成 token 返回之间的时间，主要反映图像预处理、视觉编码、prefill 和第一次 decode 的综合开销。总耗时表示完整回复生成完成所需时间。平均吞吐表示生成 token 数量与总耗时之比，decode 吞吐表示排除首 token 路径后，自回归生成阶段的平均 token 生成速率。远端计算耗时表示执行节点实际模型前向计算时间，通信非计算开销表示远程调用总耗时与远端计算耗时之差，用于衡量序列化、网络传输和协议处理成本。

\begin{table}[t]
\centering
\caption{实验环境配置}
\small
\begin{tabular}{ll}
\toprule
项目 & 配置 \\
\midrule
主控节点 & 待实测填写 \\
执行节点 & 待实测填写 \\
边缘板卡 & RK3588 / 待实测填写 \\
CPU & 待实测填写 \\
GPU & 待实测填写 \\
NPU & 待实测填写 \\
内存 & 待实测填写 \\
操作系统 & 待实测填写 \\
推理后端 & CPU/GPU/RKNN \\
模型 & Qwen3-VL-2B-Instruct / Qwen3-VL-8B-Instruct \\
网络 & 以太网，TC 限速 \\
\bottomrule
\end{tabular}
\end{table}

\subsection{端到端推理性能}

本实验比较不同执行配置下的端到端性能，包括纯软件执行、视觉 NPU 加速、文本子图 NPU 加速和分布式层执行。纯文本输入用于评估语言模型解码性能，图文输入用于评估视觉编码和多模态 prefill 对首 token 延迟的影响。

\begin{table*}[t]
\centering
\caption{不同执行配置下的端到端性能}
\scriptsize
\begin{tabularx}{\textwidth}{@{}>{\raggedright\arraybackslash}X >{\raggedright\arraybackslash}X >{\raggedright\arraybackslash}X >{\raggedright\arraybackslash}X >{\raggedright\arraybackslash}X >{\raggedright\arraybackslash}X >{\raggedright\arraybackslash}X@{}}
\toprule
模型 & 配置 & 输入 & 首 token 延迟 & 总耗时 & 平均吞吐 & Decode 吞吐 \\
\midrule
Qwen3-VL-2B & 纯软件执行 & 文本 & 待测 & 待测 & 待测 & 待测 \\
Qwen3-VL-2B & 纯软件执行 & 图文 & 待测 & 待测 & 待测 & 待测 \\
Qwen3-VL-8B & 纯软件执行 & 文本 & 待测 & 待测 & 待测 & 待测 \\
Qwen3-VL-8B & 视觉 NPU 加速 & 图文 & 待测 & 待测 & 待测 & 待测 \\
Qwen3-VL-8B & 文本子图 NPU 加速 & 文本 & 待测 & 待测 & 待测 & 待测 \\
Qwen3-VL-8B & 分布式层执行 & 文本/图文 & 待测 & 待测 & 待测 & 待测 \\
Qwen3-VL-8B & DIAL 完整框架 & 文本/图文 & 待测 & 待测 & 待测 & 待测 \\
\bottomrule
\end{tabularx}
\end{table*}

若视觉 NPU 加速显著降低图文请求首 token 延迟，说明视觉编码器是多模态请求关键瓶颈之一。若文本子图 NPU 加速提升有限，则需要进一步分析 NPU 调用开销、数据搬移成本和子图计算量之间的关系。

\subsection{分布式层划分评估}

本实验比较单机执行和不同远程层划分方案。重点观察远端计算收益是否能够抵消通信开销，以及层段聚合是否能够降低远程调用次数。

\begin{table*}[t]
\centering
\caption{分布式层划分性能}
\scriptsize
\begin{tabularx}{\textwidth}{@{}>{\raggedright\arraybackslash}X >{\raggedright\arraybackslash}X >{\raggedright\arraybackslash}X >{\raggedright\arraybackslash}X >{\raggedright\arraybackslash}X >{\raggedright\arraybackslash}X >{\raggedright\arraybackslash}X >{\raggedright\arraybackslash}X@{}}
\toprule
部署方式 & 远程层段 & 首 token 延迟 & 总耗时 & Decode 吞吐 & 远端计算耗时 & 通信非计算开销 & 远程调用次数 \\
\midrule
单机执行 & 无 & 待测 & 待测 & 待测 & 0 & 0 & 0 \\
单执行节点 & 0 至 11 层 & 待测 & 待测 & 待测 & 待测 & 待测 & 待测 \\
单执行节点 & 0 至 17 层 & 待测 & 待测 & 待测 & 待测 & 待测 & 待测 \\
双执行节点 & 0 至 11 层，12 至 23 层 & 待测 & 待测 & 待测 & 待测 & 待测 & 待测 \\
\bottomrule
\end{tabularx}
\end{table*}

如果远端节点算力较强且网络带宽充足，分布式层执行能够降低主控节点负载并提升生成速度；若通信非计算开销随远程层数增加而快速上升，则说明层划分过细或网络条件不足，需要选择更粗粒度层段或减少远程执行范围。

\subsection{通信机制消融}

本实验验证层段聚合和轻量中间特征传输对通信开销的影响。对比对象包括逐层远程执行、层段聚合执行和层段聚合加轻量传输。逐层远程执行每层触发一次通信；层段聚合执行将同一执行节点上的多个相邻层合并；层段聚合加轻量传输进一步减少层段元数据和协议处理开销。

\begin{table*}[t]
\centering
\caption{通信机制消融实验}
\scriptsize
\begin{tabularx}{\textwidth}{@{}>{\raggedright\arraybackslash}X >{\raggedright\arraybackslash}X >{\raggedright\arraybackslash}X >{\raggedright\arraybackslash}X >{\raggedright\arraybackslash}X >{\raggedright\arraybackslash}X >{\raggedright\arraybackslash}X@{}}
\toprule
请求方式 & 远程层数 & 远程调用次数 & 写入数据量 & 读取数据量 & 通信非计算开销 & 总耗时 \\
\midrule
逐层远程执行 & 待测 & 待测 & 待测 & 待测 & 待测 & 待测 \\
层段聚合执行 & 待测 & 待测 & 待测 & 待测 & 待测 & 待测 \\
层段聚合加轻量传输 & 待测 & 待测 & 待测 & 待测 & 待测 & 待测 \\
\bottomrule
\end{tabularx}
\end{table*}

理论上，逐层远程执行的请求次数与远程层数成正比，而层段聚合执行的请求次数与远程连续层段数量相关。对于 decode 阶段，请求次数下降会直接降低每个 token 的通信往返开销。轻量中间特征传输还可以减少固定元数据开销，对长序列生成更有意义。

\subsection{核心机制消融}

本实验用于验证 DIAL 三阶段框架中各个机制的增量贡献。消融从纯软件单机执行开始，逐步加入预设层段、层段聚合、轻量中间特征传输、图像尺度控制和 NPU 辅助执行。若某一机制引入后首 token 延迟、decode 吞吐或通信非计算开销发生明显变化，即可说明该机制在整体框架中的作用。

\begin{table*}[t]
\centering
\caption{DIAL 核心机制消融实验}
\scriptsize
\begin{tabularx}{\textwidth}{@{}>{\raggedright\arraybackslash}X >{\raggedright\arraybackslash}X >{\raggedright\arraybackslash}X >{\raggedright\arraybackslash}X >{\raggedright\arraybackslash}X >{\raggedright\arraybackslash}X >{\raggedright\arraybackslash}X >{\raggedright\arraybackslash}X >{\raggedright\arraybackslash}X@{}}
\toprule
配置 & 层段预设 & 层段聚合 & 轻量传输 & 图像尺度控制 & NPU 辅助 & 首 token 延迟 & Decode 吞吐 & 通信非计算开销 \\
\midrule
单机软件基线 & 否 & 否 & 否 & 否 & 否 & 待测 & 待测 & 0 \\
仅分布式层段 & 是 & 否 & 否 & 否 & 否 & 待测 & 待测 & 待测 \\
加入层段聚合 & 是 & 是 & 否 & 否 & 否 & 待测 & 待测 & 待测 \\
加入轻量传输 & 是 & 是 & 是 & 否 & 否 & 待测 & 待测 & 待测 \\
加入图像尺度控制 & 是 & 是 & 是 & 是 & 否 & 待测 & 待测 & 待测 \\
DIAL 完整框架 & 是 & 是 & 是 & 是 & 是 & 待测 & 待测 & 待测 \\
\bottomrule
\end{tabularx}
\end{table*}

该消融实验对应本文的主要创新点。层段预设主要影响单设备模型加载压力和远端计算占比；层段聚合主要影响远程调用次数和 decode 阶段逐 token 延迟；轻量传输主要影响序列化和协议处理开销；图像尺度控制主要影响图文请求的首 token 延迟；NPU 辅助执行主要影响视觉编码和规则子图的计算耗时。通过逐项加入机制，可以避免只给出完整系统性能而无法解释性能来源的问题。

\subsection{图像尺度控制评估}

图文请求中，图片尺寸影响视觉 token 数量，并进一步影响视觉编码耗时和文本 prefill 长度。本实验比较固定图像尺度和首 token 延迟约束策略下的性能差异。固定尺度用于反映传统输入预处理方式，延迟约束策略用于验证 DIAL 是否能够在满足响应时间约束的同时尽量保留视觉细节。

\begin{table*}[t]
\centering
\caption{图像尺度控制对多模态推理的影响}
\scriptsize
\begin{tabularx}{\textwidth}{@{}>{\raggedright\arraybackslash}X >{\raggedright\arraybackslash}X >{\raggedright\arraybackslash}X >{\raggedright\arraybackslash}X >{\raggedright\arraybackslash}X >{\raggedright\arraybackslash}X >{\raggedright\arraybackslash}X@{}}
\toprule
尺度策略 & 图片边长 & 视觉后端 & 首 token 延迟 & 视觉编码耗时 & Prefill 耗时 & 质量备注 \\
\midrule
固定尺度 & 336 & CPU/NPU & 待测 & 待测 & 待测 & 待测 \\
固定尺度 & 384 & CPU/NPU & 待测 & 待测 & 待测 & 待测 \\
固定尺度 & 448 & CPU/NPU & 待测 & 待测 & 待测 & 待测 \\
固定尺度 & 512 & CPU/NPU & 待测 & 待测 & 待测 & 待测 \\
延迟约束策略 & 自动选择 & CPU/NPU & 待测 & 待测 & 待测 & 待测 \\
\bottomrule
\end{tabularx}
\end{table*}

该实验用于确定边缘部署中的默认图片输入策略。若较小尺寸已能满足摔倒检测、场景问答等任务需求，则限制图片边长可以显著降低首 token 延迟；若任务需要读取细节或小目标，则需要在延迟和识别能力之间折中。延迟约束策略的目标不是固定选择最小图片，而是在给定首 token 延迟上界的条件下选择可承受的最大尺度。

\subsection{网络带宽敏感性}

分布式推理是否有效取决于远程计算收益和通信开销之间的平衡。本文通过限制主控节点与执行节点之间的带宽，评估不同网络条件下 DIAL 的性能。

\begin{table*}[t]
\centering
\caption{网络带宽对分布式推理的影响}
\scriptsize
\begin{tabularx}{\textwidth}{@{}>{\raggedright\arraybackslash}X >{\raggedright\arraybackslash}X >{\raggedright\arraybackslash}X >{\raggedright\arraybackslash}X >{\raggedright\arraybackslash}X >{\raggedright\arraybackslash}X >{\raggedright\arraybackslash}X@{}}
\toprule
带宽 & 往返时延 & 远程层段 & 首 token 延迟 & 总耗时 & 通信非计算开销 & Decode 吞吐 \\
\midrule
10 Mbps & 待测 & 待测 & 待测 & 待测 & 待测 & 待测 \\
50 Mbps & 待测 & 待测 & 待测 & 待测 & 待测 & 待测 \\
100 Mbps & 待测 & 待测 & 待测 & 待测 & 待测 & 待测 \\
1000 Mbps & 待测 & 待测 & 待测 & 待测 & 待测 & 待测 \\
\bottomrule
\end{tabularx}
\end{table*}

低带宽场景下，中间 hidden state 传输可能成为主要瓶颈，此时应减少远程切分或选择更粗粒度层段；高带宽场景下，远端节点的计算能力更容易转化为端到端收益。该实验能够界定 DIAL 在不同边缘网络条件下的适用范围。

\section{讨论}

DIAL 的主要价值在于将边缘多模态大模型部署中的资源约束、通信约束和异构后端约束统一到同一协同推理框架中。与依赖运行时任意模型切分的方法相比，DIAL 采用预设层段和层段聚合，更适合边缘硬件后端受限的实际环境。对于 RKNN 这类编译式后端，预设子图和固定输入形状更容易保证运行稳定性；对于自回归语言模型，层段聚合能够直接减少 decode 阶段的跨节点往返次数。

DIAL 的多模态设计也体现了边缘部署与模型语义一致性的平衡。视觉 token 区间对齐和 deepstack 特征融合保证了图像信息仍以模型原生方式进入语言模型；图像尺度控制则把首 token 延迟作为可优化目标，使系统能够根据应用场景选择不同的视觉细节保留程度。该机制尤其适合监控告警、巡检问答等对响应时间敏感但视觉细节要求不同的任务。

从适用范围看，DIAL 更适合模型规模超过单设备稳定承载能力、且边缘节点之间具备一定网络带宽的场景。当远端节点计算能力明显强于主控节点时，层段预设和层段聚合更容易带来端到端收益；当网络带宽较低或远端节点算力不足时，跨节点 hidden state 传输可能抵消远端计算收益，此时应减少远程层段或采用单机异构执行。对于图文请求，若任务依赖小目标、文字识别或高频细节，图像尺度控制需要设置更宽松的首 token 延迟约束；若任务以告警和粗粒度场景理解为主，则可以采用更严格的延迟约束以换取更快响应。

当前方案仍存在不足。首先，层划分仍主要依赖部署阶段 profiling，尚未根据实时网络状态和设备负载自动调整。其次，图像尺度控制需要与具体任务质量指标结合，不能仅凭延迟确定最优尺度。再次，文本 NPU 子图对模型转换、输入长度 bucket 和运行时版本敏感，不同模型规模下需要单独验证。最后，本文主要面向单请求交互式推理，对多请求并发、批处理调度和视频流连续帧复用的支持仍有进一步扩展空间。

\section{结论}

本文提出 DIAL，一种面向边缘异构平台的多模态大模型分布式协同推理框架。针对 Qwen3-VL 等视觉语言模型在边缘部署时面临的模型规模大、首 token 延迟高、单设备资源不足和硬件后端适配复杂等问题，DIAL 将协同推理抽象为“层段预设、层段聚合、异构辅助”的三阶段框架。层段预设在部署阶段确定 Transformer 连续层段的设备归属，层段聚合在推理阶段减少远端连续层的通信往返，异构辅助将视觉编码器和规则文本子图映射到 NPU 并保留软件回退路径。

在多模态输入方面，DIAL 通过视觉 token 区间对齐、主视觉 embedding 注入和 deepstack 特征融合保持 Qwen3-VL 的跨模态语义一致性，并通过首 token 延迟约束下的图像尺度控制降低图文请求关键路径耗时。本文进一步设计了端到端性能、分布式层划分、通信机制、核心机制消融、图像尺度控制和网络带宽敏感性实验，用于验证各项机制的有效性。未来工作将在现有预设拓扑机制基础上引入自动层划分、运行时自适应调度和视频流连续帧优化，以进一步提升系统在复杂边缘环境中的部署效率和推理性能。

\section*{致谢}

本文工作得到相关实验平台支持。

\begin{thebibliography}{99}

\bibitem{ref1} Y. Kang, J. Hauswald, C. Gao, A. Rovinski, T. Mudge, J. Mars, and L. Tang, “Neurosurgeon: Collaborative intelligence between the cloud and mobile edge,” ACM SIGARCH Computer Architecture News, vol. 45, no. 1, pp. 615-629, 2017.

\bibitem{ref2} Z. Hao, G. Xu, Y. Luo, H. Hu, J. An, and S. Mao, “Multi-agent collaborative inference via DNN decoupling: Intermediate feature compression and edge learning,” IEEE Transactions on Mobile Computing, vol. 22, no. 10, pp. 6041-6055, 2022.

\bibitem{ref3} M. Qin, C. Sun, J. Hofmann, and D. Vucinic, “Disco: Distributed inference with sparse communications,” in Proceedings of the IEEE/CVF Winter Conference on Applications of Computer Vision, 2024, pp. 2432-2440.

\bibitem{ref4} J. Singh, E. Fallahzadeh, B. Adams, and A. E. Hassan, “On the impact of black-box deployment strategies for edge AI on latency and model performance,” arXiv preprint arXiv:2403.17154, 2024.

\bibitem{ref5} A. Kwon, Z. Li, S. Zhuang, Y. Sheng, L. Zheng, C. H. Yu, J. E. Gonzalez, H. Zhang, and I. Stoica, “Efficient memory management for large language model serving with PagedAttention,” in Proceedings of the ACM SIGOPS Symposium on Operating Systems Principles, 2023.

\bibitem{ref6} NVIDIA, “TensorRT-LLM: Open-source library for optimizing large language model inference,” 2023.

\bibitem{ref7} Qwen Team, “Qwen-VL: A versatile vision-language model for understanding, localization, text reading, and beyond,” arXiv preprint arXiv:2308.12966, 2023.

\bibitem{ref8} Rockchip, “RKNN Toolkit and RKNN Runtime User Guide,” Rockchip Developer Documentation, 2024.

\end{thebibliography}

\end{document}
